% Actual class: 01
% Slide deck: Lecture/Chapter XX, title/version
% Exact scope: pp. XX--YY
% NTULearn supplement: video title, HH:MM:SS--HH:MM:SS, or None
% Human verified: No

\section{第 01 节课：主题}\label{lecture:SCXXXX-L01}

\lecturemeta
  {SCXXXX-L01}
  {YYYY-MM-DD}
  {Lecture/Chapter XX slides, title/version}
  {pp. XX--YY}
  {None, or video title and timestamps}
  {No}

\subsection{学习目标}

学完这部分后，应当能够：

\begin{itemize}
  \item 准确定义……；
  \item 推导或应用……；
  \item 识别……成立的假设和失效情形。
\end{itemize}

\lecturetopic{SCXXXX-L01-T01}{知识点标题}

\keyidea{用自己的话写出核心思想。}

\begin{definition}
  写出精确定义，并明确对象、符号和假设。
\end{definition}

\subsubsection{动机与机制}

解释为什么需要这一概念，以及系统或算法怎样一步步工作。

\subsubsection{推导补全}

课件若省略推导，在这里补全步骤；明确说明这是补充推导，不冒充课件原文。

\subsubsection{微型例子}

这里只保留帮助理解概念的最小例子。完整题目和完整解答移到习题区。

\subsubsection{边界、失败模式与常见误区}

\begin{itemize}
  \item 可能失效的假设：……
  \item 常见混淆：……
  \item \verify{仍需核对的结论、公式、页码或解释。}
\end{itemize}

\subsubsection{相关练习与代码作业}

\begin{itemize}
  \item \relatedexercise{SCXXXX-TUT01-Q01}{Tutorial 1, Question 1}
  \item \codeassignmentref{SCXXXX-LAB01}{assignments/scxxxx/lab01/README.md}
\end{itemize}

\subsection{外部补充}

把不属于指定课堂材料、但有助于理解的背景单独放在这里，并给出可追踪来源。

\subsection{一分钟回顾}

用三到五句话总结本次课，不引入未经说明的新结论。
